\section{Background}
\label{sec:background}

\subsection{Enumeration and symmetry}

There are $8$ magic squares of order $3$, all images of the Lo Shu under the dihedral group $D_4$ of the square (four rotations, four reflections); there are $7{,}040$ of order $4$, forming $880$ $D_4$-orbits first listed by Fr\'enicle de Bessy (published posthumously in 1693) and verified many times since \cite{beck2010enumeration}; and $275{,}305{,}224$ of order $5$, counted by Schroeppel in 1973 and reported by Gardner \cite{gardner1976}, i.e.\ $34{,}413{,}153$ orbits.
For $n\ge6$ only statistical estimates are known; see Trump's compilation \cite{trump2012enumeration}.
Since $D_4$ maps lines to lines and preserves the entry set, it maps magic squares to magic squares, and two squares in the same orbit are called \emph{essentially different} from squares in other orbits.
The $880$ order-four squares were further sorted by Dudeney into twelve groups according to the pattern of complementary pairs $(a,17-a)$ \cite{dudeney1917amusements}; we use these groups only as a classical reference partition in \Cref{sec:gem-geometry-880}.
Classical constructions (Siamese for odd $n$, diagonal-complement for $n\equiv0\pmod 4$, LUX or Strachey for $n\equiv2\pmod4$) show that magic squares exist for all $n\ge3$ \cite{andrews1960magic}; we use the Siamese squares as a family on which hull computations can be pushed to large $n$.

\subsection{Covariance and the inertia tensor}
\label{sec:bg_inertia}

For a finite point set $\mathbf p_1,\dots,\mathbf p_N\in\R^3$ with mean $\bar{\mathbf p}$, the covariance matrix is $\Sigma_{jk}=\frac1N\sum_i(p_{i,j}-\bar p_j)(p_{i,k}-\bar p_k)$.
With unit masses at the points, the inertia tensor about the centre of mass is
\begin{equation}
I_{jk}=\sum_i\bigl(\|\mathbf p_i-\bar{\mathbf p}\|^2\delta_{jk}-(p_{i,j}-\bar p_j)(p_{i,k}-\bar p_k)\bigr)
=N\bigl(\delta_{jk}\operatorname{tr}\Sigma-\Sigma_{jk}\bigr),
\label{eq:inertia_cov}
\end{equation}
where $\delta_{jk}$ is the Kronecker delta \cite{goldstein2002classical}.
Thus the off-diagonal entries of $I$ are $-N$ times the covariances and the diagonal entries are $N$ times the sum of the two other variances.
Loly's planar moment of inertia \cite{loly2004invariance} is a different object: it places the \emph{value} $s_{ij}$ as a mass at the planar point $(x_{ij},y_{ij})$ and computes $\sum s_{ij}(x_{ij}^2+y_{ij}^2)$.
The two are nevertheless related through the gem's point set: Loly's quantity is an affine function of the single degree-two cross-moment $\Cov(Z,X^2+Y^2)$ of that point set (\Cref{cor:loly}), and this is the only ingredient of the three-dimensional covariance $\Sigma$ that it uses.

\subsection{Two tools from permutation statistics}
\label{sec:bg_perm}

Let $\pi$ be a uniformly random permutation of $\{1,\dots,N\}$ and $a,b\in\R^N$ fixed vectors.
The quantity $T_a=\sum_i a_i b_{\pi(i)}$ is a \emph{linear permutation statistic}.
Its mean is $N\bar a\bar b$ and, for two such statistics with the same $b$ (Wald--Wolfowitz \cite{waldwolfowitz1944}, polarised),
\begin{equation}
\Cov(T_a,T_{a'})=\frac{1}{N-1}\sum_i(a_i-\bar a)(a'_i-\bar a')\sum_j(b_j-\bar b)^2 .
\label{eq:ww_identity}
\end{equation}
Hoeffding's combinatorial central limit theorem \cite[Thm.~3]{hoeffding1951} gives asymptotic normality of $\sum_i c_N(i,\pi(i))$ when, for every integer $r\ge3$, $N^{-1}\sum_{i,j}|d_N(i,j)|^r$ is $O\bigl((N^{-1}\sum_{i,j}d_N(i,j)^2)^{r/2}\bigr)$, $d_N$ being the doubly centred array; for product arrays $c_N(i,j)=a_ib_j$ the condition reduces to a bound on the standardised absolute moments of $a$ and $b$.
In \Cref{sec:random_energy} the energies of a random arrangement are quadratic forms in such statistics; \eqref{eq:ww_identity} gives their means exactly, and the distribution of a quadratic form in Gaussian variables is computed by Imhof's inversion formula \cite{imhof1961}.

\subsection{Polytopes of magic squares}
\label{sec:bg_polytopes}

The polytope literature on magic squares concerns a different polytope from the one studied here.
An $n\times n$ nonnegative integer matrix with all row and column sums equal to $t$ is a lattice point of the $t$-th dilate of the Birkhoff polytope of doubly stochastic matrices, and their number is a polynomial in $t$ of degree $(n-1)^2$ \cite{stanley1973magic,ehrhart1973magic,ziegler1995lectures}; adding the diagonal conditions and distinct entries leads to the inside-out polytopes of Beck and Zaslavsky \cite{beck2006insideout}, which gave the first enumeration of $4\times4$ magic squares of arbitrary magic sum \cite{beck2010enumeration}; the cones of magic squares and cubes have Hilbert-basis descriptions \cite{ahmed2003polyhedral}; and contingency tables with fixed margins are the same lattice points studied as a sample space \cite{diaconis1995rectangular}.
In all of these the magic squares are the \emph{points} of a polytope, whereas the Magic Gem attaches a polyhedron to a \emph{single} square; the two constructions share nothing beyond the word.
Spectral and matrix properties of a magic square \cite{loly2009spectra,nordgren2012properties,nordgren2014spectra}, the singular-value signatures and entropies of Cameron, Rogers and Loly \cite{cameron2013signatura}, and Loly's physical invariants---the moment of inertia \cite{loly2004invariance}, the inertia tensor and the multipole moments of a magic cube \cite{rogers2004inertia,rogers2005multipole}, and the study of Franklin squares \cite{loly2007franklin}---are closer to our viewpoint: they, too, attach a quantity to one square and ask what it reveals.
